% !TeX spellcheck = ru_RU
% !TEX root = vkr.tex

\section{Обзор литературы} \label{sec:related}

В настоящем разделе систематизируются основные результаты, относящиеся к предметной области исследования: методы параметрически эффективного дообучения больших языковых моделей, гипотеза поверхностного выравнивания, а также подходы к построению AI-агентов, способных использовать внешние инструменты. Обзор формирует теоретическую базу для последующей постановки эксперимента и интерпретации результатов.

\subsection{Методы дообучения больших языковых моделей}

В качестве базовой платформы в настоящей работе используется семейство моделей YandexGPT~\cite{yandexgpt2024docs}~--- генеративные авторегрессионные модели на основе архитектуры Transformer~\cite{vaswani2017attention, radford2019language}, разработанные компанией Яндекс и доступные в вариантах различного масштаба (YandexGPT Lite, YandexGPT Pro). Эти модели доступны через API Yandex Cloud и поддерживают дообучение на пользовательских данных.

Полное дообучение (full fine-tuning) предобученной языковой модели предполагает обновление всех её параметров на целевом наборе данных. Для моделей, содержащих миллиарды параметров, такой подход сопряжён с высокими вычислительными затратами, необходимостью хранения отдельной копии параметров для каждой задачи и риском катастрофического забывания~\cite{kirkpatrick2017overcoming}. В ответ на эти ограничения был разработан класс методов, объединённых под общим названием \textbf{Parameter-Efficient Fine-Tuning} (PEFT). Обзор данного направления представлен в работе Ding и др.~\cite{ding2023parameter}, опубликованной в Nature Machine Intelligence. Авторы систематизируют существующие подходы к параметрически эффективному дообучению и выделяют несколько основных стратегий: добавление небольших обучаемых модулей к фиксированной модели, модификация лишь малого числа существующих параметров, а также обучение дополнительных представлений на входе модели.

Метод \textbf{LoRA} (Low-Rank Adaptation), предложенный Hu и др.~\cite{hu2022lora}, параметризует обновление весов через произведение двух матриц низкого ранга, что позволяет обучать менее 1\% параметров модели при качестве, сопоставимом с полным дообучением. Формальное описание метода приведено в разделе~\ref{subsec:lora}.

Метод \textbf{Prompt Tuning}~\cite{lester2021power} добавляет обучаемые непрерывные векторы ко входу модели, оставляя её параметры замороженными. Подход \textbf{Adapter Layers}~\cite{houlsby2019parameter} встраивает небольшие обучаемые модули между слоями замороженной модели.

\subsection{Гипотеза поверхностного выравнивания}

Важным вопросом при дообучении языковых моделей является характер изменений, которые оно вносит: приобретает ли модель принципиально новые знания или лишь корректирует форму выдачи уже усвоенной при предобучении информации?

В работе Zhou и др.~\cite{zhou2023lima} была сформулирована \textbf{гипотеза поверхностного выравнивания} (Superficial Alignment Hypothesis). Авторы продемонстрировали, что модель LIMA, дообученная всего на 1000 тщательно отобранных примерах, способна генерировать ответы высокого качества, сопоставимые с результатами моделей, обученных на значительно больших наборах данных с использованием обучения с подкреплением на основе обратной связи от человека (RLHF)~\cite{ouyang2022instructgpt}. Согласно данной гипотезе, практически все фактические знания модель приобретает на этапе предобучения, тогда как дообучение с инструкциями (instruction fine-tuning) преимущественно обучает модель \textit{формату и стилю} взаимодействия с пользователем. Иными словами, выравнивание является <<поверхностным>>: оно затрагивает не глубинные знания модели, а лишь способ их представления.

Эти выводы согласуются с результатами работы Gudibande и др.~\cite{gudibande2023false}, в которой было показано, что модели, дообученные на выходах более мощных проприетарных моделей (так называемая имитация), улучшают стиль генерации, но не приобретают фактических знаний, присущих моделям-учителям. Модели-имитаторы демонстрируют улучшение качества ответов по субъективным оценкам (стиль, связность), однако на объективных бенчмарках, измеряющих фактическую корректность, прирост отсутствует~\cite{gudibande2023false}.

Для настоящего исследования гипотеза поверхностного выравнивания имеет ключевое значение. Если дообучение преимущественно корректирует стиль, а не знания, то выбор инструмента в контексте AI-агента может быть интерпретирован как элемент <<стиля>> генерации~--- устойчивый паттерн оформления ответа, который модель способна усвоить из относительно небольшого числа качественных примеров. Данная интерпретация обосновывает ожидание того, что целенаправленное дообучение на компактном, но качественном наборе данных способно значимо улучшить точность классификации инструментов.

\subsection{AI-агенты и выбор инструментов}

Понятие AI-агента на основе большой языковой модели подразумевает систему, в которой LLM выступает в роли центрального <<мыслящего>> компонента, способного планировать действия, обращаться к внешним инструментам и интерпретировать результаты их выполнения. Обзорная работа Xi и др.~\cite{xi2023rise} систематизирует архитектуры таких агентов и выделяет три ключевых компонента: модуль восприятия (perception), модуль планирования и рассуждений (brain), а также модуль действия (action), включающий взаимодействие с внешними инструментами.

Одним из наиболее влиятельных подходов к организации рассуждений агента является фреймворк \textbf{ReAct}, предложенный Yao и др.~\cite{yao2023react}. ReAct объединяет генерацию цепочки рассуждений (reasoning) и выполнение действий (acting) в чередующемся цикле <<мысль~--- действие~--- наблюдение>>. Модель последовательно формулирует рассуждение о текущей ситуации, выполняет действие (например, вызов API или поисковый запрос), получает наблюдение (результат действия) и использует его для следующего шага рассуждений. Такая архитектура позволяет агенту адаптивно выбирать инструменты на каждом шаге, исходя из контекста задачи.

В работе Schick и др.~\cite{schick2023toolformer} был предложен подход \textbf{Toolformer}, в котором языковая модель самостоятельно обучается вставлять вызовы внешних инструментов (калькулятор, поисковая система, переводчик и др.) в генерируемый текст. Обучение осуществляется в формате self-supervised: модель генерирует кандидатные вызовы инструментов, фильтрует их по критерию снижения перплексии и дообучается на отфильтрованных данных. Toolformer продемонстрировал, что модели способны научиться определять, \textit{когда} и \textit{какой} инструмент следует использовать, без явной разметки.

Работа Patil и др.~\cite{patil2023gorilla} непосредственно связана с задачей настоящего исследования. Авторы предложили модель \textbf{Gorilla}, дообученную для корректного вызова API из обширного каталога. Gorilla обучалась на наборе данных, сгенерированном с помощью Self-Instruct~\cite{wang2023selfinstruct}, и продемонстрировала способность точно генерировать вызовы API, превосходя по этому показателю модели GPT-4 без специализированного дообучения. Данный результат подтверждает гипотезу о том, что целенаправленное дообучение на задачу выбора и вызова инструментов способно существенно повысить качество классификации.

Обзорная работа Qin и др.~\cite{qin2024toollearning} систематизирует направление обучения моделей использованию инструментов (tool learning). Авторы выделяют четыре ключевых этапа взаимодействия модели с инструментами: понимание задачи, планирование использования инструментов, непосредственный вызов инструмента и обработка его результата. Особый интерес представляет выделенный авторами этап \textit{выбора инструмента} (tool selection)~--- задача классификации, в которой модель должна из доступного набора инструментов выбрать наиболее подходящий для решения текущей подзадачи. Именно эта задача является предметом настоящего исследования.

\textbf{Открытые бенчмарки tool use.} Для сравнительной оценки моделей в задачах function calling разработан ряд англоязычных бенчмарков. \textit{Berkeley Function Calling Leaderboard} (BFCL)~\cite{bfcl2024} предоставляет стандартизированные наборы задач на одиночный и множественный вызов API и параллельное исполнение; \textit{ToolBench}~\cite{qin2023toolllm} включает более 16~тыс. реальных API и используется для оценки обобщающей способности моделей; \textit{API-Bank}~\cite{li2023apibank} сосредоточен на этапах планирования и многошагового взаимодействия. Эти бенчмарки покрывают широкий спектр англоязычных API, но не содержат задач русскоязычного гео-домена, в котором функционирует исследуемый AI-агент. В настоящей работе используется собственный русскоязычный бенчмарк, отражающий специфику предметной области (разделы~\ref{subsec:benchmarks}--\ref{subsubsec:data_collection}).

\textbf{Выбор YandexGPT.} YandexGPT выбрана как модель, применяемая в промышленной эксплуатации исследуемого AI-агента; результаты дообучения данной модели напрямую переносимы в production. Сравнение с альтернативными русскоязычными LLM (GigaChat, Saiga, ruGPT) выходит за рамки настоящей работы и обозначено в направлениях дальнейших исследований (раздел~\ref{sec:conclusion}).

\subsection{Синтетическая генерация обучающих данных}

Подготовка качественных обучающих данных для дообучения языковых моделей является ресурсоёмкой задачей. Метод \textbf{Self-Instruct}, предложенный Wang и др.~\cite{wang2023selfinstruct}, предлагает подход к автоматической генерации обучающих примеров с использованием самой языковой модели. Процедура включает итеративную генерацию инструкций, входных данных и эталонных ответов с последующей фильтрацией некачественных и дублирующихся примеров. Авторы продемонстрировали, что дообучение модели на синтетически сгенерированных данных приводит к значимому улучшению качества следования инструкциям, приближаясь к результатам моделей, обученных на данных, размеченных вручную.

Результаты Self-Instruct~\cite{wang2023selfinstruct} согласуются с выводами LIMA~\cite{zhou2023lima}: если выравнивание является поверхностным и затрагивает преимущественно формат генерации, то для его достижения достаточно относительно небольшого числа качественных примеров, даже если эти примеры получены синтетическим путём. Для задачи классификации инструментов это означает, что эффективный обучающий набор может быть сформирован с помощью LLM-генерации при условии обеспечения разнообразия сценариев и корректности разметки. Модель Gorilla~\cite{patil2023gorilla} является практическим подтверждением данного подхода: её обучающий набор был получен методом Self-Instruct.

\subsection{Выводы обзора}

Проведённый обзор литературы позволяет сформулировать следующие ключевые наблюдения. Во-первых, методы параметрически эффективного дообучения, в особенности LoRA~\cite{hu2022lora}, позволяют адаптировать модели при минимальных вычислительных затратах~\cite{ding2023parameter}. Во-вторых, гипотеза поверхностного выравнивания~\cite{zhou2023lima, gudibande2023false} обосновывает ожидание того, что дообучение способно значимо улучшить паттерны генерации при небольшом объёме данных. В-третьих, работы Toolformer~\cite{schick2023toolformer} и Gorilla~\cite{patil2023gorilla} демонстрируют эффективность дообучения для задачи выбора инструментов~\cite{xi2023rise, qin2024toollearning}.

Вместе с тем, \textbf{открытым остаётся вопрос} о том, каким образом дообучение модели масштаба YandexGPT Lite влияет на качество классификации инструментов в реальной агентной системе, каков минимальный объём обучающих данных, достаточный для значимого улучшения, и в какой мере наблюдаемый эффект согласуется с гипотезой поверхностного выравнивания. Настоящая работа направлена на экспериментальное исследование именно этих вопросов.
